% ============================================================
%  Bibliografia — autor-data (NBR 6023 + ABNT similar)
% ============================================================

% --- Bibliografia ---
\chapter[Bibliografia]{Bibliografia}


\epigraph{Referências acadêmicas e fontes primárias consultadas ao longo do livro, formatadas em estilo autor-data. As citações no corpo do texto remetem ao sobrenome do autor e ao ano de publicação.}

\renewcommand{\bibsection}{\section*{}}
\renewcommand{\bibname}{}

% Estilo natbib 'plainnat' autor-data; tamanho de fonte reduzido para caber na página.
\begingroup
\small
% Estilo 'abntex2-alf' = ABNT NBR 6023 autor-data alfabético
% Configurado via \usepackage[alf]{abntex2cite} no main.tex
\bibliography{referencias}
\endgroup

% --- Fontes primárias (textos antigos citados no livro) ---
\section*{Fontes Primárias}

\begin{description}[leftmargin=2em,style=nextline]
  \item[Enoque (Livro de)] Texto apócrifo judaico, séc. III a.C. In: \textsc{Charlesworth}, James H. (Ed.). \textit{The Old Testament Pseudepigrapha}. Vol. 1. New York: Doubleday, 1983.
  \item[Apocalipse de Pedro] Texto apócrifo cristão, séc. II. In: \textsc{Elliott}, J. K. (Ed.). \textit{The Apocryphal New Testament}. Oxford: Clarendon Press, 1993.
  \item[Dante Alighieri] A Divina Comédia. Ed. crítica de Giorgio Petrocchi. Milano: Mondadori, 1966--1967.
  \item[Visão de Túndalo (Visio Tnugdali)] Ed. \textsc{Wagner}, Albrecht. Erlangen: Deichert, 1882.
  \item[Popol Vuh] Trad. \textsc{Tedlock}, Dennis. New York: Simon \& Schuster, 1996.
  \item[Tractatus de Purgatorio Sancti Patricii] Séc. XII, manuscrito medieval irlandês sobre a entrada ao purgatório através da caverna de Cruachan.
  \item[Virgílio] \textit{Eneida}, Livro VI. Composições sobre o submundo no imaginário romano. (Texto base: Mynors, R. A. B. (Ed.), Oxford: Oxford University Press, 1969.)
\end{description}

% (limpeza removida - estava gerando "paramundefined" no PDF)
